% =========================================================
% TP2 -- Automata off-lattice: bandadas de agentes autopropulsados
% AYUDA-MEMORIA DEL PRESENTADOR. No es un entregable: la consigna pide
% presentacion (b), codigo (c) e informe (d). Este documento existe para que
% quien expone tenga a mano lo que se saco de las diapositivas, que segun
% GuiaPresentaciones.pdf (1.6) deben llevar poco texto y ningun parrafo.
%
% Los numeros de diapositiva son los del pie de presentacion.pdf, que empieza
% a contar en la de Contenido (la portada no se numera).
%
%   tectonic -X compile report/ayudamemoria.tex --outdir report/build
% =========================================================

\documentclass[10pt,a4paper]{article}

\usepackage[utf8]{inputenc}
\usepackage[T1]{fontenc}
\usepackage[spanish,es-noquoting,es-noshorthands]{babel}
\usepackage[margin=2.0cm,top=1.6cm,bottom=1.8cm]{geometry}

\usepackage{amsmath}
\usepackage{amssymb}
\usepackage{siunitx}
\sisetup{output-decimal-marker={,}, per-mode=symbol}
\usepackage{xcolor}
\usepackage{enumitem}
\usepackage{tcolorbox}
\usepackage{fancyhdr}
\usepackage{needspace}

\definecolor{azul}{RGB}{23,66,110}
\definecolor{acento}{RGB}{197,90,17}
\definecolor{gris}{RGB}{110,118,128}
\definecolor{grisclaro}{RGB}{236,238,241}

\pagestyle{fancy}
\fancyhf{}
\renewcommand{\headrulewidth}{0pt}
\fancyfoot[L]{\scriptsize\color{gris} Ayuda-memoria -- Autómata off-lattice}
\fancyfoot[C]{\scriptsize\color{gris} Grupo 10 -- Comisión S2}
\fancyfoot[R]{\scriptsize\color{gris}\thepage}

\newcommand{\va}{\ensuremath{v_a}}
\newcommand{\teq}{\ensuremath{t_{\mathrm{eq}}}}
\newcommand{\Mc}{\ensuremath{M_{\mathrm{c}}}}
\newcommand{\Cmax}{\ensuremath{\mathcal{C}_{\mathrm{max}}}}

% Encabezado de cada diapositiva: numero, titulo y tiempo previsto.
\newcommand{\diapo}[3]{%
    \needspace{5\baselineskip}%
    \vspace{0.7em}%
    \noindent\colorbox{azul}{\makebox[2.1em]{\color{white}\bfseries\small #1}}%
    \colorbox{grisclaro}{\makebox[\dimexpr\linewidth-2.1em-2\fboxsep\relax][l]{%
        \color{azul}\bfseries\small\hspace{0.4em}#2\hfill
        {\color{gris}\mdseries #3}\hspace{0.4em}}}%
    \par\vspace{0.35em}}

% Rotulos de los tres tipos de nota.
\newcommand{\decir}{\textbf{\color{azul}Decir:}\;}
\newcommand{\nums}{\textbf{\color{acento}Números:}\;}
\newcommand{\preg}{\textbf{\color{gris}Si preguntan:}\;}

\newenvironment{notas}
    {\begin{itemize}[leftmargin=1.1em,itemsep=0.15em,topsep=0.15em,parsep=0pt]}
    {\end{itemize}}

\setlength{\parindent}{0pt}

\begin{document}

% =========================================================
\begin{center}
    {\LARGE\bfseries\color{azul} Ayuda-memoria del presentador}\\[0.3em]
    {\large\color{gris} Autómata off-lattice: bandadas de agentes autopropulsados}\\[0.5em]
    \textcolor{azul!30}{\rule{0.4\textwidth}{1pt}}\\[0.5em]
    {\small Grupo 10 -- Comisión S2 \enspace\textbar\enspace Trabajo Práctico 2
     \enspace\textbar\enspace exposición de \SI{13}{min}}
\end{center}
\vspace{0.4em}

\begin{tcolorbox}[colback=grisclaro,colframe=azul!40,boxrule=0.5pt,
                  arc=2pt,left=6pt,right=6pt,top=4pt,bottom=4pt]
{\small\bfseries\color{azul} Reparto del tiempo}\par\vspace{0.3em}
{\small
\begin{tabular}{@{}llr@{}}
    Introducción (3--4)      & sistema real y modelo matemático        & \SI{2}{min}\\
    Implementación (6--7)    & arquitectura y paso de simulación       & \SI{1.5}{min}\\
    Simulaciones (9--11)     & parámetros, observables, criterio de \teq & \SI{2}{min}\\
    Resultados (13--24)      & estándar, votante, clusters, tiempos    & \SI{6.5}{min}\\
    Conclusiones (26)        &                                          & \SI{1}{min}\\
\end{tabular}}
\par\vspace{0.4em}
{\small Las diapositivas \num{2}, \num{5}, \num{8}, \num{12} y \num{25} son
separadores de sección: se pasan sin hablar. Las \num{12} diapositivas de
resultados promedian \SI{30}{s}; las del votante que repiten una figura del
estándar (15, 17, 21, 23) van más rápido porque ya se explicó qué se mira, y
eso libera tiempo para las de estándar.}
\end{tcolorbox}

\vspace{0.6em}
{\small\itshape Este documento no se entrega. Reúne el texto que se sacó de las
diapositivas para cumplir con la guía de presentaciones, que pide poco texto y
sin párrafos (1.6). Las diapositivas de resultados quedaron con la figura, el
bloque de condiciones ---que la guía exige al costado (1.7)--- y a lo sumo un
par de líneas: \textbf{toda la interpretación hay que decirla en voz alta}, y
está acá abajo.}

% =========================================================
\vspace{0.8em}
{\large\bfseries\color{azul} Introducción}
\vspace{-0.2em}

\diapo{3}{Sistema real}{\SI{45}{s}}
\begin{notas}
    \item \decir El fenómeno aparece en bandadas de aves, cardúmenes y colonias
          bacterianas: agentes que se autopropulsan y ajustan su dirección a la
          de sus vecinos.
    \item \decir Lo llamativo es que la interacción es \emph{puramente local}
          ---no hay líder ni campo externo que fije una dirección--- y aun así
          el conjunto adopta espontáneamente una dirección común.
    \item \decir Hay un único parámetro de control, el ruido $\eta$, y el
          sistema exhibe una transición continua orden--desorden. Es un sistema
          fuera del equilibrio.
    \item \preg El cuadro es una corrida nuestra a $\rho = \SI{4}{m^{-2}}$,
          $\eta = \SI{0.5}{rad}$, en $t = \SI{3e3}{s}$.
\end{notas}

\diapo{4}{Modelo matemático}{\SI{75}{s}}
\begin{notas}
    \item \decir $N$ partículas puntuales, todas con la misma rapidez $v$; el
          único grado de libertad es la dirección $\theta_i$.
    \item \decir La actualización es \emph{sincrónica}: todo lo de $t+1$ se
          calcula con el estado en $t$.
    \item \decir Las dos ecuaciones de arriba son comunes a los dos escenarios.
          Lo único que cambia es $\Theta_i$, la regla de interacción.
    \item \decir El \textbf{estándar} promedia las direcciones de todo el
          entorno; el \textbf{votante} elige uno solo al azar y lo copia. Ésa es
          toda la diferencia, y produce comportamientos muy distintos.
    \item \decir La propia partícula entra en el promedio y en el sorteo, de
          modo que una partícula aislada conserva su dirección.
    \item \nums $\eta$ en radianes, $\Delta\theta_i$ uniforme en
          $[-\eta/2,\,\eta/2]$. La distancia $d$ es la de mínima imagen, por el
          contorno periódico.
    \item \preg El $\arctan$ se implementa con \texttt{atan2} para no perder el
          cuadrante.
\end{notas}

% =========================================================
\vspace{0.7em}
{\large\bfseries\color{azul} Implementación}
\vspace{-0.2em}

\diapo{6}{Arquitectura del motor}{\SI{30}{s}}
\begin{notas}
    \item \decir El motor de simulación está en C++ y el post-proceso en
          Python; se comunican por archivos de texto.
    \item \decir Ésa es la separación que pide la consigna: el motor nunca
          dibuja y el post-proceso nunca simula, así la velocidad de
          reproducción de la animación no queda atada a la de la simulación.
    \item \decir Los archivos no se detallan en la diapositiva porque la guía
          pide dejar fuera el formato de entrada/salida (2.2), pero conviene
          nombrarlos: \texttt{static.txt} lleva los parámetros de la corrida,
          \texttt{dynamic.txt} las posiciones y ángulos paso a paso, y
          \texttt{observables.txt} una fila \texttt{t va S} por paso.
    \item \preg Se separan \texttt{dynamic} y \texttt{observables} porque los
          barridos escalares no necesitan las posiciones: para las \num{5040}
          corridas se escribe sólo el segundo archivo.
\end{notas}

\diapo{7}{El paso de simulación}{\SI{60}{s}}
\begin{notas}
    \item \decir Tres pasos por iteración: buscar vecinos, calcular la nueva
          dirección, mover.
    \item \decir El detalle importante es el \textbf{buffer aparte}. Sin él, las
          partículas de índice alto se alinearían con direcciones ya
          actualizadas y la actualización dejaría de ser sincrónica.
    \item \decir El desplazamiento usa la dirección \emph{vieja}, y después se
          intercambian los buffers.
    \item \decir Los vecinos se buscan con el \textbf{CIM}: la caja se divide en
          $\Mc\times\Mc$ celdas y cada partícula se compara sólo contra su celda
          y las ocho adyacentes.
    \item \decir Para que ese vecindario alcance hace falta $L/\Mc > r_c$, lo
          que con $L = \SI{10}{m}$ y $r_c = \SI{1}{m}$ acota $\Mc \le 9$;
          usamos el máximo admisible.
    \item \decir Como control implementamos también la fuerza bruta
          $\mathcal{O}(N^2)$: con la misma semilla las dos dan trayectorias
          idénticas, así que el CIM no pierde ni duplica pares.
    \item \preg \emph{¿Es el mismo CIM del TP1?} Sí, es el mismo código sin
          modificar ---el archivo es idéntico---, junto con la geometría
          periódica. Eso es justamente lo que hace comparables los tiempos de
          la diapositiva \num{22}.
\end{notas}

% =========================================================
\vspace{0.7em}
{\large\bfseries\color{azul} Simulaciones}
\vspace{-0.2em}

\diapo{9}{Sistema particular}{\SI{35}{s}}
\begin{notas}
    \item \decir Caja de \SI{10}{m} de lado con contorno periódico, radio de
          interacción \SI{1}{m}, rapidez \SI{0.03}{m/s} y $\Delta t = \SI{1}{s}$.
    \item \decir Barremos $\eta$ de $0$ a \SI{5}{rad} con paso \SI{0.25}{rad},
          las tres densidades de la consigna y los dos modelos.
    \item \decir Agregamos un barrido \textbf{extendido} a densidades bajas
          ---$1/\pi$, $1/2\pi$ y $1/3\pi$--- porque son las únicas para las que
          $\langle k\rangle$ cae por debajo del umbral de percolación y $S$ deja
          de estar pegado a la unidad. Se usa sólo en el estudio de clusters.
    \item \nums $N = 200$, \num{400}, \num{800} en el barrido principal;
          \num{32}, \num{16}, \num{11} en el extendido. \num{20} realizaciones
          por caso (semillas \num{1} a \num{20}), \num{5040} corridas en total.
    \item \nums Las corridas de baja densidad duran \SI{2e4}{s} en vez de
          \SI{5e3}{s} porque el transitorio de agregación es mucho más largo.
    \item \preg Las semillas se reutilizan entre modelos: a igual $\rho$ y
          semilla, estándar y votante parten de la misma configuración inicial.
\end{notas}

\diapo{10}{Observables}{\SI{40}{s}}
\begin{notas}
    \item \decir La \textbf{polarización} \va{} mide cuán alineado está el
          sistema: es el módulo de la velocidad media, normalizado. Vale $1$ si
          todas van en la misma dirección.
    \item \decir No baja a cero: con direcciones independientes la suma es una
          caminata al azar, y el piso por tamaño finito es $\sqrt{\pi/4N}$. Ese
          piso reaparece en todas las curvas.
    \item \decir La \textbf{componente gigante} $S$ mide cuán conectado está:
          se arma un grafo con un nodo por partícula y una arista por cada par a
          distancia menor que $r_c$, y $S$ es la fracción de nodos en la
          componente conexa más grande.
    \item \nums Piso $\sqrt{\pi/4N}$: \num{0.0627}, \num{0.0443} y \num{0.0313}
          para $N = 200$, \num{400} y \num{800}. Medimos \num{0.064},
          \num{0.045} y \num{0.032}.
    \item \preg La componente gigante se calcula por barrido en anchura sobre
          las mismas listas de vecinos que ya devolvió el CIM: no cuesta nada
          extra.
\end{notas}

\diapo{11}{Del observable temporal al escalar}{\SI{45}{s}}
\begin{notas}
    \item \decir El valor escalar hay que medirlo sólo en el estacionario, así
          que usamos un criterio automático y uniforme, sin umbrales elegidos a
          mano.
    \item \decir Tomamos la curva \emph{promedio} de las \num{20} realizaciones,
          calculamos su media en la segunda mitad de la corrida ---que ya está
          en el estacionario--- y definimos \teq{} como el inicio del primer
          bloque de \SI{50}{s} cuya media cruza esa referencia.
    \item \decir Como el transitorio es monótono, el cruce marca su final. Y
          como se evalúa sobre el promedio, \teq{} es único por caso, no depende
          de la semilla, y es el mismo para \va{} y para $S$.
    \item \decir El escalar es el promedio de \emph{todas} las muestras con
          $t \ge \teq$ de las \num{20} realizaciones, y la barra es el desvío
          estándar de esas mismas muestras.
    \item \preg La barra mezcla la fluctuación temporal y la dispersión entre
          semillas a propósito: es la dispersión del observable en el
          estacionario, no el error de la media.
    \item \preg \emph{¿Todos los casos llegan al estacionario?} \teq{} se
          determina en los \num{126}. Como control se compara la media del
          tercer y del cuarto cuarto de la corrida: en \num{124} casos difieren
          menos que dos desvíos entre bloques; en los dos restantes
          ($\rho = 2$, $\eta = 2{,}5$ y $3{,}75$ rad, en plena transición) la
          diferencia es \num{0.03}, o sea $0{,}3\,\sigma_{\va}$, un orden menor
          que la barra.
\end{notas}

% =========================================================
\vspace{0.7em}
{\large\bfseries\color{azul} Resultados}
\vspace{-0.2em}

\diapo{13}{Animaciones características}{\SI{35}{s}}
\begin{notas}
    \item \decir Tres corridas con la misma densidad. Cada flecha es una
          partícula, con origen en su posición y color según el ángulo de su
          velocidad; la longitud es fija porque la rapidez es común a todas.
    \item \decir Izquierda, estándar con $\eta = \SI{0.5}{rad}$: color casi
          uniforme, todas van para el mismo lado. Centro, estándar con
          $\eta = \SI{4}{rad}$: colores repartidos, sin dirección común.
    \item \decir Derecha, el votante con el \emph{mismo} $\eta = \SI{0.5}{rad}$
          que ordenaba al estándar: no hay bandada única sino \textbf{dominios}
          de dirección que se arman y se disuelven. Ésa es la firma de copiar a
          uno solo en vez de promediar.
    \item \preg Las animaciones se generan con corridas aparte que sí guardan
          las posiciones; los barridos escalares no las guardan.
\end{notas}

\diapo{14}{Modelo estándar: evolución temporal}{\SI{40}{s}}
\begin{notas}
    \item \decir Tres ruidos a la misma densidad, una sola realización cada uno,
          para que se vea la fluctuación real y no la suavizada por el promedio.
    \item \decir Con $\eta = \SI{0.5}{rad}$ sube a $\va \simeq 1$ en unos
          \SI{200}{s} y prácticamente no fluctúa. Con \SI{2}{rad} hay un plateau
          intermedio con fluctuación visible. Con \SI{5}{rad} nunca ordena:
          oscila en torno al piso de tamaño finito.
    \item \decir Las verticales punteadas marcan \teq{} de cada caso; las
          horizontales, el promedio en el estacionario.
    \item \decir A la derecha, $S$ se queda pegado a $1$ en las tres: a
          $\rho = \SI{4}{m^{-2}}$ el sistema está percolado desde la condición
          inicial. Lo retomamos en clusters.
\end{notas}

\diapo{15}{Modelo de votante: evolución temporal}{\SI{25}{s}}
\begin{notas}
    \item \decir Mismos tres ruidos, misma densidad, misma semilla que la
          anterior, ahora con el votante.
    \item \decir Con $\eta = \SI{0.5}{rad}$, que ordenaba al estándar, el votante
          se estabiliza en $\va \approx 0{,}3$ y oscila entre \num{0.1} y
          \num{0.7}: la amplitud es comparable a la media. Son los dominios que
          se vieron en la animación.
    \item \decir No hay ningún mecanismo que atenúe el ruido, así que la
          fluctuación es la característica dominante del estacionario. Por eso
          la barra de error se define sobre las muestras y no sobre los
          promedios por realización.
    \item \nums $\va = \num{0.30(14)}$ en $\eta = \SI{0.5}{rad}$, contra
          \num{0.986(3)} del estándar.
\end{notas}

\diapo{16}{Modelo estándar: polarización contra ruido}{\SI{40}{s}}
\begin{notas}
    \item \decir La curva decae de forma \emph{continua} desde $\va = 1$, sin
          salto abrupto: es la transición continua que describe Vicsek.
    \item \decir A igual ruido, más densidad da más orden. La explicación es
          directa: cada partícula promedia sobre
          $\langle k\rangle = \rho\,\pi r_c^2$ vecinos, y promediar atenúa el
          ruido como $1/\sqrt{\langle k\rangle}$.
    \item \nums $\eta(\va = 0{,}5)$ pasa de \SI{2.9}{rad} a \SI{3.5}{rad} entre
          $\rho = 2$ y \SI{8}{m^{-2}}. En $\eta = \SI{3}{rad}$, $\va$ va de
          \num{0.48(12)} a \num{0.62(3)}.
    \item \preg La barra crece con el ruido porque en el desorden la
          fluctuación del observable es intrínsecamente mayor, no porque haya
          menos estadística.
\end{notas}

\diapo{17}{Modelo de votante: polarización contra ruido}{\SI{45}{s}}
\begin{notas}
    \item \decir Misma figura que la anterior, mismos ejes y los mismos tres
          colores, ahora con el votante. Conviene ir y volver una vez entre las
          dos diapositivas: el cambio se ve solo.
    \item \decir Primero, la escala del eje $\eta$: el estándar tarda hasta
          \SI{3.5}{rad} en llegar a la mitad, el votante ya cayó en
          \SI{0.2}{rad}. Un orden de magnitud antes.
    \item \decir Y sobre todo, \textbf{se dan vuelta los colores}. En el estándar
          el verde ---$\rho = \SI{8}{m^{-2}}$--- va arriba y el azul
          ---$\rho = \SI{2}{m^{-2}}$--- abajo; en el votante es exactamente al
          revés.
    \item \decir Es decir que la densidad, que en el estándar \emph{favorece} el
          orden, en el votante lo \emph{destruye}: el efecto tiene signo
          opuesto.
    \item \decir La razón es la regla de interacción. El estándar promedia sobre
          $\langle k\rangle = \rho\,\pi r_c^2$ vecinos y eso atenúa el ruido como
          $1/\sqrt{\langle k\rangle}$, así que más vecinos ayudan. El votante
          copia a uno solo elegido al azar: no promedia nada, y más vecinos sólo
          significan más direcciones distintas de donde copiar.
    \item \decir Un punto que sí comparten: en $\eta = 0$ los dos dan
          $\va = 1$. Sin ruido, copiar y promediar llevan al mismo consenso. Es
          una verificación cruzada de las dos implementaciones.
    \item \nums $\eta(\va = 0{,}5)$ a $\rho = \SI{8}{m^{-2}}$: \SI{0.2}{rad} en
          el votante contra \SI{3.5}{rad} en el estándar. A $\rho = \SI{4}{m^{-2}}$
          (los valores del informe): \SI{0.3}{rad} contra \SI{3.2}{rad}. A
          $\rho = 2$: \SI{0.4}{rad} contra \SI{2.9}{rad}.
    \item \preg \emph{¿Dónde está la comparación superpuesta?} En el informe, con
          los dos modelos en los mismos ejes a $\rho = \SI{4}{m^{-2}}$. Acá se
          comparan estas dos diapositivas, que comparten ejes y colores.
    \item \preg Las curvas del votante van con línea de trazos y marcador
          hueco, y las del estándar con línea llena y marcador lleno, para que
          no se confundan si se muestran juntas.
\end{notas}

\diapo{18}{Clusters: animaciones a baja densidad}{\SI{25}{s}}
\begin{notas}
    \item \decir Con las tres densidades de la consigna el grafo de vecinos está
          percolado siempre y $S$ no aporta información: por eso mostramos los
          clusters a $\rho = 1/\pi$, donde hay en promedio un vecino por
          partícula y $S$ recorre todo su rango.
    \item \decir A la izquierda, con poco ruido, todo terminó agregado en un
          único grupo. A la derecha, con $\eta = \SI{5}{rad}$, queda
          fragmentado en grupos de dos o tres.
    \item \preg \emph{¿Por qué densidades que no pide la consigna?} Porque con
          $\rho = 2$, \num{4} y \num{8} el ítem de clusters queda sin contenido:
          $S \ge 0{,}90$ en todo el barrido. Las tres pedidas están todas
          reportadas; las bajas se agregan para que $S$ tenga rango.
\end{notas}

\diapo{19}{Clusters: evolución temporal de la componente gigante}{\SI{40}{s}}
\begin{notas}
    \item \decir Ésta es la figura del ítem (d): $S(t)$ para las tres densidades
          de la consigna, más las dos bajas, todas al mismo ruido
          $\eta = \SI{2}{rad}$ y promediadas sobre las \num{20} realizaciones.
    \item \decir Las tres de la consigna ---verde, roja, azul--- están en
          $S \simeq 1$ desde el primer paso: el sistema nace percolado. Sus
          corridas duran \SI{5e3}{s}, por eso terminan antes.
    \item \decir Las dos bajas ---violeta y rosa--- son otra cosa: $S$ parte de
          \num{0.15} y sube por agregación. Dos partículas que se encuentran se
          alinean, y al alinearse dejan de separarse. La unión sólo se deshace
          por el ruido.
    \item \decir El transitorio es diez veces más largo: \teq{} de \SI{3.65e3}{s}
          y \SI{3.25e3}{s} contra \SI{200}{}--\SI{950}{s}. Por eso estas
          corridas duran \SI{2e4}{s}.
    \item \decir En el panel izquierdo se ve además que a baja densidad \va{}
          sube \emph{junto con} $S$: adelanta el acople de la diapositiva 22.
    \item \nums \teq{} en $\eta = 2$: \num{200}, \num{650}, \num{950},
          \num{3650}, \SI{3250}{s} para $\rho = 8$, \num{4}, \num{2}, $1/\pi$,
          $1/3\pi$. En otros puntos del barrido bajo llega a \SI{5.6e3}{s}.
    \item \preg En $\eta = 0$ a baja densidad \teq{} es mucho mayor (hasta
          \SI{1.96e4}{s}): sin ruido la agregación no se detiene hasta que todo
          el sistema es un único grupo, y la curva sigue subiendo hasta el final.
          El valor escalar ahí es $S \simeq 1$ y $\va \simeq 1$, que es robusto
          aunque la ventana estacionaria sea corta.
\end{notas}

\diapo{20}{Modelo estándar: componente gigante contra ruido}{\SI{35}{s}}
\begin{notas}
    \item \decir $\rho = 8$ y $2$ quedan pegadas a $S = 1$: sus
          $\langle k\rangle$ valen \num{25} y \num{6.3}, muy por encima del
          umbral de percolación de un grafo geométrico aleatorio, que está en
          torno a \num{4.5}.
    \item \decir $1/\pi$ y $1/3\pi$ sí decaen de forma monótona, hasta
          $S \simeq 0{,}2$.
    \item \decir Lo más interesante: en $\eta = 0$ percolan igual, con
          $\langle k\rangle$ de \num{1} y \num{0.33}. No contradice el umbral,
          porque ése vale para un grafo \emph{aleatorio} y sin ruido las
          posiciones dejan de serlo: la agregación es irreversible.
    \item \decir O sea que \textbf{lo que fragmenta es el ruido, no la
          densidad}.
    \item \nums En $\eta = 0$: $S = \num{1.000(1)}$ y \num{0.991(40)}. Ruido de
          media componente: \SI{2.74}{rad} y \SI{2.12}{rad}. En $\eta = 5$,
          $S \approx 0{,}2$ son grupos de \num{6.4} y \num{2.5} partículas.
\end{notas}

\diapo{21}{Modelo de votante: componente gigante contra ruido}{\SI{20}{s}}
\begin{notas}
    \item \decir Mismos ejes y colores. El votante reproduce el mismo cuadro
          ---altas pegadas a 1, bajas que decaen, $\eta = 0$ percolado--- pero se
          fragmenta antes.
    \item \nums Ruido de media componente a $\rho = 1/\pi$: \SI{1.19}{rad}
          contra \SI{2.74}{rad} del estándar. A $1/3\pi$: \SI{1.13}{rad}.
    \item \decir La razón es la misma de la polarización: copiar a un único
          vecino no atenúa el ruido, así que los grupos se deshacen con menos.
\end{notas}

\diapo{22}{Modelo estándar: polarización contra componente gigante}{\SI{40}{s}}
\begin{notas}
    \item \decir Cada punto es un valor de $\eta$ del barrido, con las cuatro
          densidades. La figura separa dos regímenes de manera nítida.
    \item \decir A $\rho = 8$ y \SI{2}{m^{-2}} los puntos se apilan contra el
          borde derecho: \va{} recorre casi todo $[0,1]$ mientras $S$ se queda
          cerca de la unidad. Alineación y conectividad están \emph{desacopladas}
          y \va{} no es función de $S$.
    \item \decir A $1/\pi$ y $1/3\pi$ la relación es monótona creciente y las
          dos densidades caen casi sobre la misma curva: la alineación no puede
          propagarse más allá del grupo conexo, así que $S$ fija cuánta
          polarización global es alcanzable.
    \item \nums Pearson \num{0.45} y \num{0.61} contra \num{0.985} y \num{0.994}.
    \item \preg La degeneración es explícita: para $S \approx 0{,}957$ el
          estándar da $\va = \num{0.68}$ en $\eta = \SI{2.25}{rad}$ y
          \num{0.11} en \SI{4.5}{rad}. Misma conectividad, polarizaciones muy
          distintas.
    \item \preg El corrimiento vertical entre $1/\pi$ y $1/3\pi$ es de tamaño
          finito: en el extremo desordenado cada una satura en su propio piso
          $\sqrt{\pi/4N}$ (\num{0.157} y \num{0.267}).
\end{notas}

\diapo{23}{Modelo de votante: polarización contra componente gigante}{\SI{20}{s}}
\begin{notas}
    \item \decir Mismos ejes y colores; los mismos dos regímenes.
    \item \decir A $\rho = \SI{8}{m^{-2}}$ el Pearson es incluso \emph{negativo},
          $-0{,}48$: confirma que la correlación residual de las densidades altas
          no tiene significado físico.
    \item \nums A baja densidad el acople se mantiene: \num{0.968} y
          \num{0.987}.
    \item \decir Conclusión de las dos: el desacople es un efecto de los
          parámetros que fija la consigna, no del modelo.
\end{notas}

\diapo{24}{Tiempos de ejecución del CIM}{\SI{35}{s}}
\begin{notas}
    \item \decir El cronómetro abarca sólo la búsqueda de vecinos ---incluido el
          armado de las celdas--- y excluye la dinámica y la escritura a disco.
    \item \decir Para comparar con el TP1 hay que igualar la geometría, así que
          estas corridas replican su caja: $L = \SI{20}{m}$ y $\Mc = 13$, no las
          del resto del trabajo.
    \item \decir La fuerza bruta crece como $\mathcal{O}(N^2)$; el CIM, con
          exponente \num{1.34}. No es \num{1} porque $L$ es fijo: al subir $N$
          sube la densidad y con ella el número de candidatos por celda.
    \item \decir Lo que el CIM reduce no es el orden sino el \textbf{prefactor}:
          cada partícula se compara contra el \SI{5.3}{\percent} del sistema en
          lugar de contra todo él.
    \item \decir Para $N$ chico la fuerza bruta gana, porque el costo fijo de
          armar las celdas supera al de comparar unas pocas decenas de pares.
    \item \decir La comparación con el TP1 es directa porque el algoritmo es el
          mismo código: las dos curvas tienen la misma forma y la misma
          pendiente, y el corrimiento vertical es del equipo, no del algoritmo.
    \item \nums Bruta: de \SI{0.035}{\milli\second} a \SI{5.25}{\milli\second}
          entre $N = 113$ y \num{1150}. CIM: de \SI{0.025}{\milli\second} a
          \SI{0.53}{\milli\second}. Cruce en $N \approx 90$; factor \num{10} en
          $N = 1150$. Corrimiento respecto del TP1: entre \num{1.1} y \num{1.7}.
    \item \preg \emph{¿Por qué no se ve la mejora de orden?} Porque el CIM no
          cambia el orden a densidad creciente; lo que cambia es la constante.
          Con $\rho$ fija y $N$ creciendo sí sería lineal.
    \item \preg Las mediciones son muy sensibles a la carga de la máquina: al
          reejecutar el binario del TP1 sin modificar sobre un equipo con otros
          procesos, $N = 1150$ pasó de \SI{0.904}{\milli\second} a
          \SI{3.756}{\milli\second}. Ambos conjuntos se tomaron con el equipo
          descargado.
\end{notas}

% =========================================================
\vspace{0.7em}
{\large\bfseries\color{azul} Conclusiones}
\vspace{-0.2em}

\diapo{26}{Conclusiones}{\SI{60}{s}}
\begin{notas}
    \item \decir El modelo estándar exhibe una transición continua
          orden--desorden controlada por $\eta$, y la densidad favorece el
          orden.
    \item \decir El votante es un orden de magnitud más sensible al ruido y la
          densidad le juega con signo opuesto; en $\eta = 0$ los dos coinciden
          en $\va = 1$.
    \item \decir Con las densidades de la consigna la componente gigante no es
          un grado de libertad: $S \ge 0.90$ siempre. Al bajar la densidad ese
          desacople desaparece y $S$ queda fuertemente acoplado a \va{}. O sea
          que el desacople es un efecto de los parámetros, no del modelo.
    \item \decir Sin ruido el sistema percola aun con un vecino por partícula:
          lo que fragmenta es el ruido.
    \item \decir Y el CIM es \num{10} veces más rápido que la fuerza bruta en
          $N = 1150$, con el mismo escalamiento que medimos en el TP1.
\end{notas}

\vspace{0.9em}
\begin{tcolorbox}[colback=grisclaro,colframe=gris!45,boxrule=0.5pt,
                  arc=2pt,left=6pt,right=6pt,top=4pt,bottom=4pt]
{\small\bfseries\color{azul} Preguntas probables del final}\par\vspace{0.3em}
{\small
\begin{itemize}[leftmargin=1.1em,itemsep=0.2em,topsep=0.2em,parsep=0pt]
    \item \emph{¿Por qué $\Mc = 9$ y no más?} Porque el CIM sólo mira las ocho
          celdas adyacentes, y para que ese vecindario cubra el radio de
          interacción hace falta $L/\Mc > r_c$. Con $L = \SI{10}{m}$ y
          $r_c = \SI{1}{m}$ eso da $\Mc \le 9$, y conviene el máximo porque
          celdas más chicas dejan menos candidatos por comparar.
    \item \emph{¿Por qué la actualización sincrónica?} Es lo que define el
          autómata celular: si se actualizara en el lugar, el resultado
          dependería del orden de recorrido de los índices.
    \item \emph{¿Cómo saben que el CIM está bien?} Contra la fuerza bruta: con
          la misma semilla las dos producen trayectorias idénticas, así que no
          pierde ni duplica pares.
    \item \emph{¿Por qué agregaron densidades que no pide la consigna?} Porque
          con $\rho = 2$, \num{4} y \num{8} el grafo está percolado desde la
          condición inicial y el ítem de clusters queda sin contenido: $S \ge
          0.90$ en todo el barrido. Las tres densidades pedidas están todas
          reportadas; las bajas se agregan para que $S$ tenga rango.
    \item \emph{¿Por qué \va{} no baja a cero?} Por el tamaño finito: con
          direcciones independientes la suma de $N$ vectores es una caminata al
          azar y su módulo medio da $\sqrt{\pi/4N}$.
\end{itemize}}
\end{tcolorbox}

\end{document}
